% ============================================================
\section{行为批判：公理体系的系统性失效}\label{sec:behavioral}
% ============================================================

理性选择框架（完备 + 传递 + 效用最大化）是分析基准，
但大量受控实验与现场数据表明：人的实际选择\textbf{系统性地}偏离公理。
批判分两类：
\textbf{情境依赖}（偏好随选择集、表述方式、社会情境改变——偏好并非预先存在）
与\textbf{非深思选择}（选择不经由两两比较与最大化，而由直觉、启发式驱动）。
每条批判都精确对应某条公理或某条结构性质的失效。

\subsection{批判一：情境依赖——诱饵效应}

\begin{example}[诱饵效应]\label{ex:decoy}
两属性商品空间：质量（纵轴）与价格（横轴，左好右差）。
\begin{itemize}
  \item 选项 $A$：高质量、高价格；
  \item 选项 $B$：低质量、低价格；
  \item 诱饵 $C$：质量略低于 $A$、价格高于 $A$（被 $A$ 严格支配）。
\end{itemize}
选择集 $\{A,B\}$ 中，$A$ 与 $B$ 各获约一半选择；
加入\textbf{无人选择的}诱饵 $C$ 后，选 $A$ 的比例显著上升。
\end{example}

\begin{center}
\begin{tikzpicture}[scale=1.1]
  \draw[axisline] (0,0) -- (6.4,0) node[below left, axislabel] {价格（向右递增）};
  \draw[axisline] (0,0) -- (0,4.6) node[below left, axislabel] {质量};
  \node[nodeblue, label=above left:{$A$}] at (3.4,3.4) {};
  \node[nodeblue, label=above left:{$B$}] at (1.2,1.1) {};
  \node[nodeblue, label=below right:{$C$}] at (4.6,3.0) {};
  % dominance region of A over C
  \draw[dashed, MidBlue] (3.4,3.4) -- (4.6,3.4) -- (4.6,3.0);
  \node[align=left] at (3.35,4.25) {$A$：高质高价，$B$：低质低价，互不占优};
  \node[align=left] at (4.45,1.9) {$C$（诱饵）：质量低于 $A$、\\ 价格高于 $A$，被 $A$ 严格支配};
\end{tikzpicture}
\end{center}
\figexplan{横轴价格向右递增、纵轴质量向上递增。$A$ 高质高价、$B$ 低质低价，二者互不占优；诱饵 $C$ 质量低于 $A$ 且价格高于 $A$，因而被 $A$ 严格支配（虚线折线标出支配关系）。$C$ 本身无人选择，却显著提高 $A$ 的选择份额，直接违反理性框架的正则性。}

\begin{proposition}[诱饵效应违反正则性]\label{prop:decoy}
记 $P(x,S)$ 为从选择集 $S$ 中选中 $x$ 的（群体）频率。
理性框架的\textbf{正则性}要求：$T\subseteq S\Rightarrow P(x,T)\ge P(x,S)$——
选择集扩大不会提升原有元素被选中的份额。
诱饵实验中 $P\bigl(A,\{A,B,C\}\bigr)>P\bigl(A,\{A,B\}\bigr)$，
直接违反正则性；在确定性选择函数语言下，
这对应于违反显示偏好理论中的\textbf{无关选项独立性}（IIA）。
\end{proposition}

\begin{remark}[机制]
诱饵提供了\textbf{局部比较基准}：$A$ 对 $C$ 的明显优势被归因为 $A$ 的“品质”，
而 $A$ 与 $B$ 之间需要跨属性权衡（困难比较）。
选择依赖“哪个比较容易论证”，而非稳定的内在排序——
偏好在选择时刻被\textbf{构造}，而非被\textbf{揭示}。
\end{remark}

\subsection{批判二：框架效应}

\begin{example}[生存率与死亡率]\label{ex:framing}
同一项治疗方案的两种表述：
\begin{itemize}
  \item 表述一：“治疗后存活率为 $90\%$”；
  \item 表述二：“治疗后死亡率为 $10\%$”。
\end{itemize}
两种表述在信息上\textbf{完全等价}，但表述一下接受率显著更高。
类似地，“含脂 $25\%$”与“瘦肉 $75\%$”的同一食品获得不同评价。
\end{example}

\begin{remark}[违反的公理：描述不变性]
理性框架隐含\textbf{描述不变性}：选择只依赖选项的结果结构，
与表述方式无关。框架效应表明描述本身进入偏好——
同一结果集在不同“参照点 + 损益编码”下被评估为不同方案。
这直接动摇效用表示的前提：$u$ 的定义域不再是“客观结果”，
而是“被框架化后的结果”，$\succeq$ 在客观结果集上甚至\textbf{无法良定义}。
\end{remark}

\subsection{批判三：偏好反转}

\begin{example}[定价与选择的矛盾]\label{ex:reversal}
两个不确定性选项：$P$ 型（高概率、小收益）与 $\$$ 型（低概率、大收益）。
受试者在\textbf{直接选择}中选 $P$ 型，
但对两者的\textbf{定价}（最低愿售价）却给 $\$$ 型定更高价。
若偏好稳定，$P\succ\$$ 意味着 $P$ 的定价应更高——观察到的是系统性反转。
\end{example}

\begin{remark}[违反的公理：程序不变性]
选择与定价是同一潜在偏好的两种\textbf{导出方式}，程序不变性要求它们一致。
反转表明：不同引出程序激活不同的认知过程
（选择任务突出概率维度、定价任务突出金额维度），
产出不同的“偏好”——不存在一个同时解释两者的稳定 $\succeq$。
\end{remark}

\subsection{批判四：非深思选择——启发式、默认与现状偏差}

\begin{example}[启发式选择]
购买复杂商品时，以“选最知名品牌”“选销量最高”“选中间价位”
代替逐属性比较——这些规则\textbf{不遍历}全部选项，
不满足“对所有备选作两两比较”的完备性操作。
\end{example}

\begin{example}[默认效应与现状偏差]\label{ex:default}
同一保险计划，设为\textbf{默认加入}（需主动退出）时参保率，
远高于设为\textbf{默认退出}（需主动加入）时。
若偏好稳定，默认设置（一个不改变任何选项可行性的“框”）不应改变选择。
现状偏差的形式化：存在\textbf{参照依赖}——
以现状 $r$ 为参照点时 $\succeq_r$，以 $r'$ 为参照点时 $\succeq_{r'}$，
二者不必一致；特别地，“失去”的权重大于等量“得到”（损失厌恶），
使无差异曲线在参照点处出现\textbf{折点}。
\end{example}

\begin{center}
\begin{tikzpicture}[scale=1.05]
  \draw[axisline] (0,0) -- (6.2,0) node[below left, axislabel] {商品 1};
  \draw[axisline] (0,0) -- (0,4.4) node[below left, axislabel] {商品 2};
  \node[nodeblue, label=below right:{参照点 $r$}] at (3.0,2.0) {};
  % kinked IC through r
  % 折点无差异曲线：左支陡（损失域）、右支平（收益域），两支均向右下倾斜
  \draw[thick, MidBlue] (0.4,3.72) -- (3.0,2.0) -- (5.9,1.18);
  \node[figlabel] at (1.05,3.98) {无差异曲线};
  % 虚线：无损失厌恶时的直线基准（过 r，斜率与损失域相同）
  \draw[dashed, MidBlue] (0.4,2.74) -- (5.5,0.37);
  \node[align=left,figlabel] at (1.30,1.80) {损失域（左上）：斜率陡峭，\\放弃一单位需更大补偿};
  \node[align=left,figlabel] at (4.70,2.60) {收益域（右下）：斜率平缓，\\同等收益只需较小付出};
\end{tikzpicture}
\end{center}
\figexplan{无差异曲线在参照点 $r$ 处折曲，且两支均向右下倾斜：左支（损失域，失去商品 1、获得商品 2）斜率陡，右支（收益域）斜率平缓。虚线是“若无损失厌恶”时的直线基准，其斜率与损失域相同；实际曲线在 $r$ 左侧高于该基准（更陡）、在右侧也高于该基准（更平）。两支斜率之比即损失厌恶程度：放弃一单位所需的补偿大于获得同量带来的收益。}

\begin{remark}[禀赋效应：买卖价差]
同一物品，拥有者的最低愿售价（WTA）系统性高于非拥有者的最高愿购价（WTP）。
在 §3.7 拟线性框架下，补偿变化与等值变化应\textbf{相等}（无财富效应时）；
即使允许财富效应，观察到的价差也远超收入效应可解释的幅度。
参照依赖（$r$=“已拥有”对 $r'$=“未拥有”）给出统一解释：
“卖出”被编码为损失，“买入”被编码为得到。
\end{remark}

\subsection{批判五：锚定与偏好构建}

\begin{example}[锚定]
先让受试者写下证件号码后两位，再问其愿意为某商品支付多少：
号码大者出价系统性更高——一个\textbf{明显无关}的数字影响了估值。
\end{example}

\begin{remark}
锚定是“偏好构建论”的极端证据：估值过程从最近激活的数字出发做不充分调整，
说明许多“偏好”没有预先存储的数值，而是在被询问的当下\textbf{临时构造}。
这与理性框架的基本设定（偏好是先于选择的稳定序）正面冲突。
\end{remark}

\subsection{效应与公理的对照}

\begin{center}
\small
\begin{tabular}{lll}
\toprule
行为效应 & 被违反的公理 / 性质 & 失效机制 \\
\midrule
诱饵效应 & 正则性 / IIA & 局部比较基准改变相对吸引力 \\
折中效应 & 正则性 / IIA & 选项因“居中”而获益 \\
框架效应 & 描述不变性 & 损益编码随表述改变 \\
偏好反转 & 程序不变性 & 不同引出程序激活不同维度权重 \\
默认 / 现状偏差 & 偏好稳定性 & 参照点进入偏好本身 \\
禀赋效应（WTA$>$WTP） & 补偿对称性（§3.7） & 损失权重 $>$ 得到权重 \\
锚定 & 偏好预先存在性 & 估值当下构造、不充分调整 \\
冲动选择 & 深思比较假设 & 启发式 / 直觉替代权衡 \\
\bottomrule
\end{tabular}
\end{center}

\begin{remark}[文献线索]
本章各效应的奠基性实验与文献：诱饵（吸引）效应源于 Huber, Payne \& Puto (1982) 的“占优违背”实验，并经 Tversky \& Simonson (1993) 系统化为 attraction effect；框架效应由 Tversky \& Kahneman (1981) 在“亚洲疾病问题”中确立；偏好反转由 Lichtenstein \& Slovic (1971) 发现、Grether \& Plott (1979) 以受控实验严格复核；默认 / 现状偏差的经典现场证据见 Johnson \& Goldstein (2003) 对欧洲器官捐献制度的跨国比较，Thaler \& Benartzi (2004) 的“明天储蓄更多”（Save More Tomorrow）计划则是其政策应用；禀赋效应由 Thaler (1980) 提出、Kahneman, Knetsch \& Thaler (1990) 以“咖啡杯实验”给出干净识别；锚定效应见 Tversky \& Kahneman (1974)，Ariely, Loewenstein \& Prelec (2003) 用“随机数字轮盘”排除了“理性锚定”的替代解释；上述效应的统一理论框架是 Kahneman \& Tversky (1979) 的\textbf{前景理论}（参照依赖 + 损失厌恶 + 概率权重）。
\end{remark}

\subsection{综合：理性基准与行为视角}

\begin{center}
\small
\begin{tabular}{lll}
\toprule
 & 理性选择模型 & 行为视角 \\
\midrule
偏好 & 稳定、预先存在 & 选择时刻被构造 \\
决策规则 & 两两比较 + 效用最大化 & 启发式、直觉、局部比较 \\
对框架的敏感性 & 无（描述不变） & 系统性敏感 \\
参照点 & 无 & 核心（损失厌恶） \\
财富效应 & 拟线性情形为零 & 普遍存在且随情境变化 \\
理论角色 & 基准与近似 & 对偏离方向的解释与预测 \\
\bottomrule
\end{tabular}
\end{center}

\begin{remark}[如何对待批判]
批判并未使理性框架失效，而是标定了它的\textbf{适用边界}：
(i) 作为\textbf{基准}——偏离是有方向的（诱饵使目标选项份额上升），
理性模型给出“无偏离”的零点，行为模型给出偏离项；
(ii) 作为\textbf{近似}—— stakes 高、学习机会多、激励强的环境中，
选择更接近公理；一次性、低激励、高复杂度的选择中偏离最大；
(iii) 公理的可检验化（第二章的显示偏好路线）使“偏离”可以被度量：
对每个选择数据集，可计算其违反 SARP 的程度（不一致指数），
把“理性程度”从是非判断变为连续变量。
\end{remark}

\begin{remark}[本章收束]
从公理到表示（§3.1--§3.4），从表示到结构（§3.5--§3.7），再到结构的边界（§3.8）：
效用函数是\textbf{序信息的压缩编码}，其存在性由可数序稠密条件精确刻画，
其特殊形式（可加、拟线性）由相应的序条件精确刻画，
而其失效处（情境、框架、参照点）同样可以在公理语言中精确指认。
这套“条件 $\leftrightarrow$ 表示”的方法论，是后续消费者理论、
不确定情形下效用理论（期望效用及其批判）的直接模板。
\end{remark}
